\documentclass{nature}
\usepackage{graphicx}
\usepackage{amsmath}

\title{Title of the paper in one sentence}

\author{First Author$^{1}$, Second Author$^{1,2}$ \& Third Author$^{2}$}

\begin{document}

\maketitle

\begin{affiliations}
\item Department, University, City, Country.
\item Institute, City, Country.
\end{affiliations}

\begin{abstract}
A summary paragraph of about 150 words without references: the background,
the question, the main result and what it means.
\end{abstract}

The main text has no headings in a Letter and few in an Article. Introduce the
question and how earlier work left it open\cite{watson1953}.

The main result follows. Equations are numbered like this:
\begin{equation}
  E = mc^2,
  \label{eq:main}
\end{equation}
and figures are referred to as Fig.~\ref{fig:overview}\cite{jumper2021}.

\begin{figure}
  \centering
  \includegraphics[width=0.6\linewidth]{example-image}
  \caption{\textbf{A one-sentence title for the figure.} Then what each panel
  shows. Upload an image and use its file name here.}
  \label{fig:overview}
\end{figure}

Close with what the result means and what comes next.

\begin{methods}
\subsection{Samples.}
Describe the materials and methods in enough detail to reproduce the work.

\subsection{Statistics.}
Describe the statistical tests and how the sample sizes were chosen.
\end{methods}

\bibliographystyle{naturemag}
\bibliography{references}

\begin{addendum}
\item[Acknowledgements] Thank the people and funding that supported this work.
\item[Author contributions] F.A. and S.A. designed the study; T.A. wrote the paper.
\item[Competing interests] The authors declare no competing interests.
\item[Correspondence] Requests for materials should be addressed to F.A.\\ (first@example.org).
\end{addendum}

\end{document}
